\documentclass[10pt]{beamer}
\usetheme{Madrid}
\usepackage[T1]{fontenc}
\usepackage[utf8]{inputenc}
\usepackage[brazil]{babel}
\usepackage{lmodern}
\usepackage{hyperref}
\usepackage{listings}
\usepackage{xcolor}
\setbeamertemplate{navigation symbols}{}

\lstset{
  language=C++,
  basicstyle=\ttfamily\small,
  keywordstyle=\color{blue},
  commentstyle=\color{gray},
  stringstyle=\color{red},
  breaklines=true,
  showstringspaces=false,
}

\title{Estrutura de Programas: Funções e Structs}
\subtitle{Curso de Microcontroladores}
\author{}
\date{}

\begin{document}

\begin{frame}
  \titlepage
\end{frame}

\begin{frame}{O desafio de um programa grande}
  \begin{itemize}
    \item Quando seu programa tem centenas de linhas, fica muito difícil de ler, manter e corrigir.
    \item Programas de robótica, como um seguidor de linha, precisam:
    \begin{itemize}
      \item Ler sensores constantemente.
      \item Fazer cálculos matemáticos complexos (como PID).
      \item Controlar motores em tempo real.
      \item Talvez comunicar via serial.
    \end{itemize}
    \item Fazer tudo isso em um único bloco de código seria complexo e geraria um código difícil de manter.
    \item A solução é \textbf{quebrar o programa em partes menores e testáveis}.
  \end{itemize}
\end{frame}

\begin{frame}{O Padrão IPO (Entrada, Processamento, Saída)}
  \begin{itemize}
    \item Todo programa segue logicamente um padrão simples: \textbf{IPO}.
    \item \textbf{Entrada}: O robô coleta dados do mundo físico (sensores).
    \item \textbf{Processamento}: O código analisa esses dados e toma decisões.
    \item \textbf{Saída}: O robô atua no mundo (motores, LEDs, etc).
    \item No caso do seguidor de linha:
    \begin{itemize}
      \item \textbf{Entrada}: Lê a linha com sensores ópticos.
      \item \textbf{Processamento}: Calcula o erro de posição e aplica PID.
      \item \textbf{Saída}: Ajusta a velocidade dos motores.
    \end{itemize}
  \end{itemize}
\end{frame}

\begin{frame}{O Padrão IPO (Entrada, Processamento, Saída)}
  \begin{itemize}
    \item O padrão IPO não representa só o que o seguidor é capaz de fazer por inteiro, com os sensores sendo entrada, o PID sendo processamento e os motores sendo saída.
    \item Ele também pode ser aplicado dentro de cada coisa, ou seja, o sensor pode ser seu IPO, sendo a entrada a luz refletida, o processamento a conversão analógica-digital e a saída o valor 0-1023.
    \item O PID também tem seu próprio IPO, onde a entrada é o erro, o processamento é o cálculo dos termos P, I e D, e a saída é a correção.
    \item O controle dos motores também tem seu IPO, onde a entrada é a correção, o processamento é a validação e ajuste da velocidade, e a saída é o sinal PWM aplicado aos motores.
    \item Ou seja, o padrão IPO é uma forma de organizar o pensamento e o código em níveis hierárquicos, desde o programa inteiro até cada função individual.
    \item Isso torna tudo muito mais simples
  \end{itemize}
\end{frame}

\begin{frame}[fragile]{\texttt{loop()}}
\begin{lstlisting}
void setup() {
  // Executado APENAS uma vez ao ligar
  Serial.begin(9600);
  pinMode(SENSOR_PIN, INPUT);
}

void loop() {
  // Executado REPETIDAMENTE, centenas de vezes por segundo
  int leitura = analogRead(SENSOR_PIN);
  Serial.println(leitura);
  delay(10);  // pequena pausa
}
\end{lstlisting}

  \begin{itemize}
    \item `setup()` roda uma única vez: você configura pinos, inicia comunicação serial, etc.
    \item `loop()` roda infinitamente
    \item Quanto mais rápido o `loop()` rodar, mais responsivo o robô será.
  \end{itemize}
\end{frame}

\begin{frame}{Por que usar funções?}
  \begin{itemize}
    \item Sim, você \textit{poderia} escrever tudo no `loop()`, mas seria pouco eficiente.
    \item Funções permitem que você \textbf{quebre o problema em partes menores}.
    \item Cada função possui uma responsabilidade única.
    \item Você testa cada função isoladamente antes de integrar.
    \item Exemplo: uma função apenas para ler sensores, outra apenas para calcular, outra para atuar.
    \item Isso também facilita trabalhar em equipe: um desenvolve `lerSensores()`, outro desenvolve `calcularPID()`.
  \end{itemize}
\end{frame}

\begin{frame}[fragile]{Exemplo: Um programa sem funções}
\begin{lstlisting}
void loop() {
  // Ler 8 sensores manualmente
  int s0 = analogRead(A0);
  int s1 = analogRead(A1);
  // ... mais 6 linhas ...
  int s7 = analogRead(A7);
  
  // Calcular média e erro (muitas linhas)
  int somatorio = s0 + s1 + ... + s7;
  int media = somatorio / 8;
  int erro = media - 500;
  
  // Controlar motores manualmente
  analogWrite(MOTOR_L, 200 - erro);
  analogWrite(MOTOR_R, 200 + erro);
}
\end{lstlisting}

  \begin{itemize}
    \item Tudo misturado no `loop()`.
    \item Se houver um bug, é difícil saber se é na leitura, no cálculo ou na saída.
    \item Código muito repetitivo (muito `analogRead()` seguido).
  \end{itemize}
\end{frame}

\begin{frame}[fragile]{O que é uma função?}
\begin{lstlisting}
int somar(int a, int b) {
  return a + b;
}

void setup() {
  int resultado = somar(5, 3);  // resultado = 8
}
\end{lstlisting}

  \begin{itemize}
    \item Função é um bloco de código reutilizável com um nome específico.
    \item `int somar(int a, int b)` significa: recebe dois inteiros e \textbf{retorna} um inteiro.
    \item `void` significa que a função não retorna nada (apenas executa).
    \item Parâmetros (`a` e `b`) são dados que você passa para a função.
    \item `return` envia um valor de volta ao código que chamou a função.
  \end{itemize}
\end{frame}

\begin{frame}[fragile]{Mini IPO: Função de Entrada}
\begin{lstlisting}
int lerSensores() {
  // ENTRADA: coleta dados
  int valores[8];
  for (int i = 0; i < 8; i++) {
    valores[i] = analogRead(sensorPinos[i]);
  }
  
  // PROCESSAMENTO: calcula a posição
  int somatorio = 0;
  for (int i = 0; i < 8; i++) {
    somatorio += valores[i];
  }
  int posicao = somatorio / 8;
  
  // SAÍDA (da função): retorna o resultado
  return posicao;
}
\end{lstlisting}

  \begin{itemize}
    \item Cada função pode ter seu próprio \textbf{mini IPO}.
    \item Esta função: lê sensores, processa, retorna um valor.
    \item Você pode testar `lerSensores()` isoladamente para garantir que funciona.
  \end{itemize}
\end{frame}

\begin{frame}[fragile]{Mini IPO: Função de Processamento}
\begin{lstlisting}
int calcularPID(int posicaoAtual) {
  // ENTRADA: recebe a posição
  static int erroAnterior = 0;
  int erro = posicaoAtual - 500;
  
  // PROCESSAMENTO: cálculo PID
  float P = erro * 0.5;
  float I = /* acumulador */ ;
  float D = (erro - erroAnterior) * 0.2;
  int correcao = (P + I + D);
  
  erroAnterior = erro;
  
  // SAÍDA: retorna a correção
  return correcao;
}
\end{lstlisting}

  \begin{itemize}
    \item Esta função recebe um valor, processa com PID, retorna a correção.
    \item Você pode testar com valores conhecidos: se posição é 500 (centro), deve retornar 0.
  \end{itemize}
\end{frame}

\begin{frame}[fragile]{Mini IPO: Função de Saída}
\begin{lstlisting}
void controlarMotores(int velocidadeEsq, 
                      int velocidadeDir) {
  // ENTRADA: recebe as velocidades (parâmetros)
  
  // PROCESSAMENTO: validação
  velocidadeEsq = constrain(velocidadeEsq, 0, 255);
  velocidadeDir = constrain(velocidadeDir, 0, 255);
  
  // SAÍDA: atua nos motores
  analogWrite(MOTOR_L, velocidadeEsq);
  analogWrite(MOTOR_R, velocidadeDir);
}
\end{lstlisting}

  \begin{itemize}
    \item `void` significa que não retorna nada, apenas executa.
    \item Recebe velocidades, valida, aplica aos pinos PWM.
    \item `constrain()` garante que os valores fiquem entre 0-255.
  \end{itemize}
\end{frame}

\begin{frame}[fragile]{O \texttt{loop()} como Orquestrador}
\begin{lstlisting}
void loop() {
  // Mini IPO 1: Ler sensores
  int posicao = lerSensores();
  
  // Mini IPO 2: Calcular correção
  int correcao = calcularPID(posicao);
  
  // Mini IPO 3: Atuar nos motores
  int velEsq = 150 - correcao;
  int velDir = 150 + correcao;
  controlarMotores(velEsq, velDir);
}
\end{lstlisting}

  \begin{itemize}
    \item O `loop()` não contém lógica complexa, apenas chamadas de função.
    \item Fica muito legível: você vê exatamente o que o robô faz em cada ciclo.
    \item Se há um problema, você sabe exatamente qual função investigar.
    \item Cada função roda rapidamente, mantendo a taxa de amostragem alta.
  \end{itemize}
\end{frame}

\begin{frame}{Estruturas: Agrupando Dados Relacionados}
  \begin{itemize}
    \item Um robô tem muitos componentes: sensores, motores, LED, bateria, etc.
    \item Em vez de criar variáveis soltas (`motorEsqVel`, `motorDirVel`, `motorEsqPino`...), usamos `struct`.
    \item `struct` é um tipo de dado personalizado que agrupa valores relacionados.
    \item Exemplo: um motor possui pino PWM, pino de direção, velocidade atual. Todos esses dados relacionados podem estar em uma `struct`.
  \end{itemize}
\end{frame}

\begin{frame}[fragile]{Definindo e Usando uma Estrutura}
\begin{lstlisting}
struct Motor {
  int pinoPWM;
  int pinoDir;
  int velocidade;
};

void setup() {
  Motor motorEsq;
  motorEsq.pinoPWM = 5;
  motorEsq.pinoDir = 6;
  motorEsq.velocidade = 0;
  
  pinMode(motorEsq.pinoPWM, OUTPUT);
  pinMode(motorEsq.pinoDir, OUTPUT);
}
\end{lstlisting}

  \begin{itemize}
    \item Você acessa propriedades com o `.` (ponto).
    \item Deixa claro que essas três variáveis pertencem ao mesmo motor.
    \item Muito mais legível do que `motor1_PWM`, `motor1_DIR`, etc.
  \end{itemize}
\end{frame}

\begin{frame}[fragile]{Funções que recebem estruturas}
\begin{lstlisting}
void controlarMotor(Motor motor, int velocidade) {
  motor.velocidade = velocidade;
  analogWrite(motor.pinoPWM, velocidade);
  
  if (velocidade > 0) {
    digitalWrite(motor.pinoDir, HIGH);  // frente
  } else {
    digitalWrite(motor.pinoDir, LOW);   // trás
  }
}

void loop() {
  controlarMotor(motorEsq, 150);
  controlarMotor(motorDir, 150);
}
\end{lstlisting}

  \begin{itemize}
    \item Você passa uma `struct` inteira para a função.
    \item A função trabalha com todos os dados do motor de uma vez.
    \item Muito mais elegante do que passar cada propriedade separadamente.
  \end{itemize}
\end{frame}

\begin{frame}[fragile]{Exemplo Completo: Seguidor Estruturado}
\begin{lstlisting}
struct Sensor {
  int pino;
  int valor;
};

struct Motor {
  int pinoPWM;
  int pinoDir;
};

Sensor sensores[8];
Motor motorEsq, motorDir;

void setup() {
  // Inicializa sensores e motores
  for (int i = 0; i < 8; i++) {
    sensores[i].pino = A0 + i;
  }
  motorEsq.pinoPWM = 5;
  motorEsq.pinoDir = 6;
  // ... etc
}
\end{lstlisting}

  \begin{itemize}
    \item Você pode ter arrays de estruturas (`sensores[8]`).
    \item Muito mais organizado do que dezenas de variáveis soltas.
  \end{itemize}
\end{frame}

\begin{frame}{Testando Cada Função Isoladamente}
  \begin{itemize}
    \item Quando você quebra o programa em funções, pode testar cada uma separadamente.
    \item \textbf{Teste de entrada}: rode apenas `lerSensores()`, coloque Serial.println() e veja se os valores fazem sentido.
    \item \textbf{Teste de processamento}: passe valores conhecidos para `calcularPID()` e verifique se o resultado está correto.
    \item \textbf{Teste de saída}: rode `controlarMotores()` e ouça o motor para garantir que a velocidade está mudando.
    \item Isso torna muito mais fácil encontrar e corrigir bugs.
    \item Você descobrir que um motor não está funcionando é bem diferente de descobrir que o cálculo do PID está errado.
  \end{itemize}
\end{frame}

\begin{frame}{Laços (Loops) dentro de Funções}
  \begin{itemize}
    \item Usar laços `for` e `while` dentro de funções é importante para controlar fluxo e repetição.
    \item Laços no `loop()` do Arduino vs laços dentro de funções:
    \begin{itemize}
      \item Laços dentro de funções têm \textbf{escopo local}, o que significa que as variáveis do laço não afetam o resto do programa.
      \item Você tem mais controle: pode quebrar, continuar, ou fazer múltiplas iterações de forma clara.
    \end{itemize}
    \item Exemplo: usar `for` para ler todos os 8 sensores dentro de `lerSensores()` é mais claro do que ter 8 linhas soltas no `loop()`.
  \end{itemize}
\end{frame}

\begin{frame}[fragile]{Cuidado com \texttt{delay()}}
\begin{lstlisting}
// ❌ RUIM: delay() no loop() bloqueia tudo
void loop() {
  lerSensores();
  delay(100);  // 100ms sem fazer nada!
  calcularPID();
}

// ✓ BOM: usar millis() para não bloquear
unsigned long tempoAnterior = 0;
void loop() {
  unsigned long tempoAtual = millis();
  if (tempoAtual - tempoAnterior >= 100) {
    tempoAnterior = tempoAtual;
    lerSensores();
    calcularPID();
  }
}
\end{lstlisting}

  \begin{itemize}
    \item `delay()` congela todo o programa. O robô não responde a sensores durante esse tempo!
    \item Para tarefas periódicas, use `millis()` e verificação de tempo.
    \item Isso mantém o robô responsivo e permite múltiplas tarefas em ``paralelo''.
  \end{itemize}
\end{frame}

\begin{frame}{Modularização: Além de Funções}
  \begin{itemize}
    \item À medida que o código cresce (centenas ou milhares de linhas), organizar tudo em um único arquivo `.ino` fica impossível.
    \item A solução é \textbf{separar o código em múltiplos arquivos}, cada um com uma responsabilidade.
    \item Exemplo de um seguidor de linha modularizado:
    \begin{itemize}
      \item `config.h` → constantes, pinos, parâmetros globais.
      \item `motor.h` / `motor.cpp` → tudo sobre controle de motores.
      \item `sensores.h` / `sensores.cpp` → leitura e processamento de sensores.
      \item `pid.h` / `pid.cpp` → cálculo do algoritmo PID.
      \item `bluetooth.h` / `bluetooth.cpp` → comunicação sem fio.
    \end{itemize}
    \item Cada módulo é \textbf{independente} e pode ser testado isoladamente.
  \end{itemize}
\end{frame}

\begin{frame}[fragile]{Arquivos `.h` (Header) - Declarações}
\begin{lstlisting}
// motor.h
#ifndef MOTOR_H
#define MOTOR_H

#include <Arduino.h>

void inicializarMotor();
void pararMotores();
void controlarMotor(int velEsq, int velDir);

#endif
\end{lstlisting}

  \begin{itemize}
    \item Arquivo `.h` contém \textbf{declarações} e \textbf{protótipos} de funções.
    \item `#ifndef` / `#define` / `#endif` previnem inclusão duplicada (include guards).
    \item É o contrato: aqui está o que o módulo oferece.
    \item Quem quer usar motor, inclui `motor.h` e pronto.
  \end{itemize}
\end{frame}

\begin{frame}[fragile]{Arquivos `.cpp` - Implementação}
\begin{lstlisting}
// motor.cpp
#include "motor.h"
#include "config.h"

void inicializarMotor() {
  for (int i = 0; i < 4; i++) {
    pinMode(MOTOR_PINS[i], OUTPUT);
  }
  pararMotores();
}

void pararMotores() {
  analogWrite(PWM_LEFT, 0);
  analogWrite(PWM_RIGHT, 0);
}

void controlarMotor(int velEsq, int velDir) {
  analogWrite(PWM_LEFT, velEsq);
  analogWrite(PWM_RIGHT, velDir);
}
\end{lstlisting}

  \begin{itemize}
    \item Arquivo `.cpp` contém a \textbf{implementação} (o código real).
    \item Inclui o `.h` correspondente (para ter acesso às declarações).
    \item Pode incluir outros módulos (`config.h`, `sensores.h`, etc.) conforme necessário.
  \end{itemize}
\end{frame}

\begin{frame}[fragile]{Organização em Pastas}
  Estrutura típica de um projeto modularizado:
\begin{lstlisting}
meu_seguidor/
├── meu_seguidor.ino (setup() e loop())
└── src/
    ├── config/
    │   └── config.h (constantes e pinos)
    ├── motor/
    │   ├── motor.h
    │   └── motor.cpp
    ├── sensores/
    │   ├── sensores.h
    │   └── sensores.cpp
    ├── pid/
    │   ├── pid.h
    │   └── pid.cpp
    └── bluetooth/
        ├── bluetooth.h
        └── bluetooth.cpp
\end{lstlisting}

  \begin{itemize}
    \item Cada módulo em sua própria pasta.
    \item `config.h` centraliza \textbf{todas} as constantes (pinos, velocidades, parâmetros).
    \item Muito fácil de encontrar e editar um componente específico.
  \end{itemize}
\end{frame}

\begin{frame}[fragile]{O arquivo `config.h`}
\begin{lstlisting}
// config.h - Centraliza TODAS as constantes
#ifndef CONFIG_H
#define CONFIG_H

// Pinos
const int LED_PIN = 22;
const int PWM_LEFT = 21;
const int PWM_RIGHT = 2;
const int MOTOR_PINS[] = {18, 19, 5, 4};

// Parâmetros PID
float KP = 0.5;
float KD = 0.4;

// Velocidade
int VEL_MAX = 80;

// Estados globais
bool start = false;
bool linhaPerdida = false;

// Timing
unsigned long lastLoopTime = 0;
const unsigned long loopInterval = 30;

#endif
\end{lstlisting}

  \begin{itemize}
    \item Todos os módulos incluem `config.h` para acessar constantes.
    \item Se você quer mudar um pino, muda em um lugar só.
    \item Evita replicação de código (DRY: Don't Repeat Yourself).
  \end{itemize}
\end{frame}

\begin{frame}[fragile]{Incluindo Módulos no `.ino`}
\begin{lstlisting}
// meu_seguidor.ino
#include "src/config/config.h"
#include "src/motor/motor.h"
#include "src/sensores/sensores.h"
#include "src/pid/pid.h"
#include "src/bluetooth/bluetooth.h"

void setup() {
  Serial.begin(115200);
  inicializarMotor();
  inicializarSensores();
  inicializarBluetooth();
}

void loop() {
  processarComandosBT();
  
  if (millis() - lastLoopTime < loopInterval) return;
  lastLoopTime = millis();
  
  lerSensoresQTR();
  if (!verificarLinhaPerdida() && start) {
    calcularPID();
  } else {
    pararMotores();
  }
}
\end{lstlisting}

  \begin{itemize}
    \item Arquivo `.ino` inclui todos os módulos que precisa.
    \item `setup()` e `loop()` ficam simples e legíveis.
    \item Cada função vem de seu respectivo módulo.
  \end{itemize}
\end{frame}

\begin{frame}{Benefícios da Modularização}
  \begin{itemize}
    \item \textbf{Legibilidade}: código organizado e fácil de navegar.
    \item \textbf{Manutenibilidade}: problema em motor? Vai direto em `motor.cpp`.
    \item \textbf{Reutilização}: use `motor.h` e `motor.cpp` em outro projeto.
    \item \textbf{Testes}: teste cada módulo isoladamente antes de integrar.
    \item \textbf{Trabalho em equipe}: um desenvolve `motor.cpp`, outro `sensores.cpp`, sem conflitos.
    \item \textbf{Escalabilidade}: adicione novos módulos (LED, buzzer, LoRa) sem mexer no código existente.
    \item \textbf{Performance}: o compilador só inclui o código que você realmente usa.
  \end{itemize}
\end{frame}

\begin{frame}[fragile]{Anti-padrão: Sem Modularização}
\begin{lstlisting}
// ❌ RUIM: tudo em um arquivo de 1000+ linhas
void setup() {
  // ... 50 linhas de inicialização ...
}

void loop() {
  // ... 200 linhas de código misturado ...
  // sensores, cálculos, motor, BT... tudo aqui!
}

// ... 50 funções espalhadas sem organização ...
\end{lstlisting}

  Consequências:
  \begin{itemize}
    \item Difícil encontrar uma função específica.
    \item Um bug pode estar em qualquer lugar.
    \item Impossível reutilizar código em outro projeto.
    \item Trabalhar em equipe é muito difícil.
  \end{itemize}
\end{frame}

\begin{frame}{Resumo: Estrutura de um Bom Programa}
  \begin{enumerate}
    \item \textbf{Defina estruturas} para agrupar dados relacionados (Motor, Sensor).
    \item \textbf{Quebre em funções}: cada função faz uma coisa bem.
    \item \textbf{Use mini IPOs}: cada função tem entrada → processamento → saída.
    \item \textbf{Mantenha o \texttt{loop()} limpo}: apenas chamadas de função.
    \item \textbf{Teste isoladamente}: valide cada função antes de integrar.
    \item \textbf{Evite \texttt{delay()}}: use `millis()` para tarefas periódicas.
    \item \textbf{Modularize com arquivos}: separe em `.h` / `.cpp` e organize em pastas.
  \end{enumerate}
  
  \begin{itemize}
    \item Um programa estruturado assim fica fácil de entender, manter e depurar.
    \item Você consegue trabalhar em equipe sem conflitos.
    \item Adicionar novos recursos é trivial.
    \item Seu código fica profissional e reutilizável.
  \end{itemize}
\end{frame}

\end{document}
